\section{从人类偏好学习奖励}

\subsection{偏好比较}

直接让人类给绝对评分困难且噪声大。更好的方式：展示两个选项，让人类选择哪个更好。

$$p(y_1 \succ y_2 \mid x) = \sigma\left(R_\phi(x, y_1) - R_\phi(x, y_2)\right)$$

\begin{importantbox}{Bradley-Terry 模型}
给定人类偏好数据 $(y_w, y_l)$（$y_w$ 被偏好），训练奖励模型 $R_\phi$ 最大化：
$$\mathcal{L}(\phi) = \mathbb{E}\left[\log \sigma\left(R_\phi(x, y_w) - R_\phi(x, y_l)\right)\right]$$

这就是 RLHF 中\textbf{奖励模型训练}的核心公式。
\end{importantbox}

\subsection{RLHF Pipeline}

\begin{enumerate}
    \item 收集人类偏好比较数据
    \item 训练奖励模型 $R_\phi$
    \item 用 RL（如 PPO）优化策略以最大化 $R_\phi$ 的输出
\end{enumerate}

\begin{warningbox}{奖励模型也会被利用}
与 goal classifier 类似，如果 RL 优化过度，策略可能找到奖励模型的漏洞而不是真正满足人类偏好。因此通常加入 KL 约束，限制策略不能偏离初始策略太远。
\end{warningbox}

\subsection{本章小结}

从偏好比较中学习奖励模型是 RLHF 的基础。Bradley-Terry 模型将偏好比较转化为奖励差的 sigmoid，训练简单高效。

